
\begin{abstract}
Operating room (OR) scheduling often prioritizes operational efficiency over patient preferences. In our study, we present a patient-centered scheduling framework developed for the orthopedic department at Hôpital Lapeyronie, the University Hospital of Montpellier, France. Formulated as a mixed-integer linear program (MILP), the model handles strict resource constraints with patient availabilities. This framework generates a set of optimized time slots, giving the patients the ability to select their preferred surgery date. We used an iterative heuristic that proposes a patient at a time, where the physician sets with the patient the best schedule date between three options, allowing the excluded ones to be considered for the other patients. The results using use case data show that the MILP model increases the theoretical capacity of the operating rooms utilization by 23\%. Benchmarks were also generated to simulate daily uses and showed that, on average, more than 2.5 slots are proposed to each patient. This tool shows that efficiency can be achieved with humane consideration by integrating operational objectives with patient satisfaction.

\end{abstract}
